\documentclass[11pt,twocolumn]{article}

\usepackage[a4paper,margin=0.75in,columnsep=0.3in]{geometry}
\usepackage[T1]{fontenc}
\usepackage{amsmath,amssymb,amsthm,mathtools}
\usepackage{bm}
\usepackage{graphicx}
\usepackage{booktabs}
\usepackage{array}
\usepackage{enumitem}
\usepackage{float}
\usepackage[numbers,sort&compress]{natbib}
\usepackage{hyperref}
\hypersetup{colorlinks=true,linkcolor=blue,citecolor=blue,urlcolor=blue}

\newtheorem{theorem}{Theorem}
\newtheorem{definition}{Definition}
\newtheorem{postulate}{Postulate}
\newtheorem{corollary}{Corollary}
\newtheorem{proposition}{Proposition}

\title{Unified Dynamics II: Dark Matter, Baryogenesis, the Cosmological Constant, and Black Hole Information from the Tetrahedral Framing Phase Transition}
\author{%
Benjamin Walker Mayes\thanks{*Contact author: ben.w.mayes@gmail.com}\\[4pt]
{\small Unaffiliated, Cedar Park, Texas 78613, USA}\\[8pt]
Arthur Eugene Mayes\\
{\small Loyola University School of Law, USA}\\[4pt]
Dustin Huff\\
{\small Unaffiliated, USA}\\[4pt]
Brayden Ross Sanders\\
{\small Unaffiliated, USA}\\[4pt]
Jeffrey A.\ Dewey\\
{\small Unaffiliated, USA}\\[4pt]
Brandon Bevin\\
{\small Georgia Southern University, USA}\\[4pt]
A.\ Christopher Beckingham\\
{\small Queen's University, Canada}\\[4pt]
Jonathan Engle\\
{\small Unaffiliated, USA}%
}
\date{}

\begin{document}
\maketitle

\begin{abstract}
Building on the transport--closure framework of Unified Dynamics (UD), we show that the $\mathbb{Z}_2$ framing phase transition on the tetrahedral IFS lattice simultaneously addresses four outstanding problems in fundamental physics.
(1)~Dark matter is identified with stable neutral composite packet relics that survived the framing transition: they carry burden density $\rho$ and interact gravitationally via the emergent metric but carry no gauge charge.
(2)~Baryogenesis arises from CP-violating dynamics during the transition: the nonperturbative framing-ordering process generates a baryon asymmetry whose pre-washout magnitude is $\sim 10^{-5}$; standard sphaleron dilution and washout are expected to reduce this to the observed $\eta_B\sim 10^{-10}$.
(3)~The cosmological constant is the residual vacuum energy of trapped framing defects: $\rho_\Lambda = n_{\mathrm{def}}\cdot\nu_0\cdot\Delta E_{\mathrm{flip}}$, with the IFS screening suppressing the naive QFT estimate by a factor $(\Delta/N)^{n_{\mathrm{IFS}}}$.
(4)~The black hole information paradox is resolved by carrier conservation: the five-component carrier $\mathcal{C}=(C_\rho,C_\theta,C_\chi,C_\sigma,C_b)$ is conserved through all transport processes including gravitational collapse; information is encoded in carrier correlations that leak through the horizon via the entanglement mechanism.
All four results follow from a single event---the electroweak-scale ordering of the $\mathbb{Z}_2$ framing field---using only constants derived in the companion paper.
\end{abstract}

%═══════════════════════════════════════════════════════════
\section{Introduction}
\label{sec:intro}
%═══════════════════════════════════════════════════════════

Four of the most significant open problems in fundamental physics---the identity of dark matter, the origin of the matter--antimatter asymmetry, the smallness of the cosmological constant, and the fate of information in black holes---are typically treated as independent puzzles requiring separate mechanisms.
We show here that within the Unified Dynamics (UD) framework~\cite{mayes2026ud,mayes2026grav}, all four problems trace to consequences of a single event: the $\mathbb{Z}_2$ framing phase transition on the tetrahedral IFS lattice.

The companion paper~\cite{mayes2026ud} established that the Standard Model gauge group $\mathrm{SU}(3)\times\mathrm{SU}(2)\times\mathrm{U}(1)$ emerges from the representation theory of $S_4$ on $K_4$, with electroweak symmetry breaking driven by the ferromagnetic ordering of the $\mathbb{Z}_2$ framing field $\sigma_i\in\{+1,-1\}$.
The gravitational branch~\cite{mayes2026grav} showed that Einstein dynamics emerges as the relaxed coarse isotropic limit of a four-channel accommodation architecture.

In the present paper, we develop the cosmological consequences of the framing transition.
When the early-universe temperature drops below the framing critical temperature $T_c$, the disordered (paramagnetic) phase orders into the ferromagnetic vacuum $\sigma_i=+1$.
This transition necessarily produces:
stable neutral relics (Sec.~\ref{sec:dm}),
a matter--antimatter asymmetry (Sec.~\ref{sec:baryo}),
residual vacuum energy (Sec.~\ref{sec:cc}),
and a mechanism for information conservation through gravitational collapse (Sec.~\ref{sec:bh}).

We use only constants derived in Ref.~\cite{mayes2026ud}: $\Delta=7$, $N=192$, $x=7/192$, $S_0=43/7$, $H=101/6$, $\kappa_C=3/(4\pi)$, $\eta_O=1/2$, $\nu_0\approx 83$~keV, and the five-component carrier $\mathcal{C}=(C_\rho,C_\theta,C_\chi,C_\sigma,C_b)$.

%═══════════════════════════════════════════════════════════
\section{The Framing Phase Transition}
\label{sec:transition}
%═══════════════════════════════════════════════════════════

The framing sector of the packet action is a ferromagnetic Ising model on the tetrahedral IFS lattice:
\begin{equation}
H_\sigma = -J_\sigma\sum_{\langle ij\rangle}\sigma_i\sigma_j,\qquad J_\sigma = \frac{1}{2},
\label{eq:Hsigma}
\end{equation}
with coordination number $z=E=6$ (edges per vertex on $K_4$).
The mean-field critical temperature is
\begin{equation}
T_c = \frac{z J_\sigma}{k_B} = \frac{6\cdot(1/2)}{k_B} = \frac{3}{k_B}.
\label{eq:Tc}
\end{equation}
In natural units with $\nu_0$ as the energy scale, $T_c = 3\nu_0\approx 249$~keV.
This is the UD-native framing transition scale.
The full electroweak symmetry-breaking process is a two-stage event:

\textbf{Stage~1} ($T\sim v_{\mathrm{EW}}\sim 100$~GeV): the $\mathrm{SU}(2)$ gauge sector acquires mass through the Higgs mechanism.
Sphaleron processes are active at this stage, satisfying the conditions for baryogenesis (Sec.~\ref{sec:baryo}).

\textbf{Stage~2} ($T\sim 3\nu_0\approx 249$~keV): the $\mathbb{Z}_2$ framing field orders ferromagnetically.
Framing defects that survive this transition become dark matter relics (Sec.~\ref{sec:dm}), and trapped defects contribute to the cosmological constant (Sec.~\ref{sec:cc}).

The structural analysis below depends on the existence of both stages; the precise relationship between $\nu_0$ and $v_{\mathrm{EW}}$ is set by RG running, a computation that remains open.

Above $T_c$, the framing field is disordered: $\langle\sigma_i\rangle=0$.
The full $\mathrm{SU}(2)\times\mathrm{U}(1)$ symmetry is unbroken.
Below $T_c$, the system orders into $\langle\sigma_i\rangle=+1$ (or $-1$; the two are related by $\mathbb{Z}_2$), breaking $\mathrm{SU}(2)\to\mathrm{U}(1)_{\tau_3}$.

The order parameter is the magnetization
\begin{equation}
m = \langle\sigma_i\rangle = \tanh\!\left(\frac{zJ_\sigma\,m}{T}\right),
\label{eq:magnetization}
\end{equation}
which rises continuously from $m=0$ at $T_c$ to $m\to 1$ at $T\to 0$.
This is a second-order phase transition with critical exponents determined by the Ising universality class in the IFS fractal dimension $d_f\approx 2$ (the Sierpi\'nski tetrahedron has Hausdorff dimension $\ln 4/\ln 2=2$).

The transition produces three types of relics:
\begin{enumerate}[label=(\roman*),leftmargin=2em,itemsep=2pt]
\item neutral packets that never acquired framing charge,
\item trapped framing defects (domain walls, vortices) that failed to annihilate,
\item a net asymmetry in $\sigma=+1$ vs $\sigma=-1$ domains from the CP-violating dynamics of the ordering process.
\end{enumerate}

%═══════════════════════════════════════════════════════════
\section{Dark Matter from Neutral Packet Relics}
\label{sec:dm}
%═══════════════════════════════════════════════════════════

\begin{theorem}[Dark matter candidate]
\label{thm:dm}
Stable neutral packets---composite tetrahedral closures of $V=4$ sites satisfying $\sum_{i\in\mathrm{composite}}\sigma_i=0$ and closure state $\chi_i=1$ for all constituent sites---interact gravitationally through their total burden density but carry no net gauge charge.
They are the UD dark matter candidate.
\end{theorem}

\textit{Proof sketch.}
A neutral packet has state $X_i=(\rho_i,\theta_i,1,0,b_i)$ where the framing $\sigma_i$ averages to zero over the packet's internal structure.
The packet action Eq.~(2) of Ref.~\cite{mayes2026ud} couples $\rho_i$ to the curvature term via $\kappa_C$, so neutral packets source the emergent metric and feel gravity.
However, gauge couplings arise from the framing ($\sigma$) and closure ($\chi$) sectors: $\mathrm{SU}(2)$ from edge-pair framing, $\mathrm{U}(1)$ from vertex phase.
A packet with $\langle\sigma\rangle=0$ has no net weak isospin and no net hypercharge---it is a gauge singlet.
It interacts only gravitationally.

The neutral seed threshold is $S_0=43/7$; the charged seed threshold is $2S_0=86/7$.
At the framing transition, the fraction of packets that remain neutral is determined by the overlap deficit:
\begin{equation}
f_{\mathrm{neutral}} = \frac{\Delta}{N} = \frac{7}{192}.
\label{eq:fneutral}
\end{equation}
The fraction that become charged (acquire definite framing) is $1-x=185/192$.

\begin{proposition}[Relic abundance ratio]
The dark-to-baryonic matter ratio is
\begin{equation}
\frac{\Omega_{\mathrm{DM}}}{\Omega_b} = \frac{S_0}{2S_0-S_0}\cdot\frac{f_{\mathrm{neutral}}}{1-f_{\mathrm{neutral}}}\cdot\frac{m_{\mathrm{neutral}}}{m_{\mathrm{baryon}}}.
\label{eq:ratio}
\end{equation}
\end{proposition}

The first factor $S_0/S_0=1$ reflects equal threshold costs.
The second factor is $x/(1-x)=7/185\approx 0.038$.
The third factor is the mass ratio, which depends on the neutral packet mass.

For the simplest estimate: if neutral packets have mass $m_n=\nu_0\cdot S_0\approx 0.51$~MeV (same as the electron), then:
\begin{multline}
\frac{\Omega_{\mathrm{DM}}}{\Omega_b}\approx\frac{7}{185}\cdot\frac{S_0}{m_p/\nu_0}\\
\approx\frac{7}{185}\cdot\frac{43/7}{938/0.083}\approx 2\times 10^{-5}.
\end{multline}
This is too small by a factor of $\sim 3\times 10^5$ compared to the observed $\Omega_{\mathrm{DM}}/\Omega_b\approx 5.3$.

The resolution: neutral packets are NOT single-site objects.
They are COMPOSITE closures---tetrahedral clusters of 4 packets forming a minimal 3D closure with $F=4$ faces, each face a triangular color singlet.
The composite neutral packet mass is
\begin{equation}
m_{\mathrm{DM}} = V\cdot\nu_0\cdot S_0 = 4\cdot 83\text{ keV}\cdot\frac{43}{7}\approx 2.04\text{ MeV}.
\end{equation}

More importantly, the relic abundance is determined by the freeze-out dynamics at $T_c$, not by the simple fraction $x$.
The Boltzmann suppression at freeze-out gives:
\begin{equation}
\frac{n_{\mathrm{DM}}}{n_b}\sim\left(\frac{m_{\mathrm{DM}}}{T_c}\right)^{3/2}e^{-m_{\mathrm{DM}}/T_c}\cdot\frac{g_{\mathrm{DM}}}{g_b},
\end{equation}
where $g_{\mathrm{DM}}=1$ (singlet) and $g_b$ counts baryonic degrees of freedom.

The key structural prediction: the dark matter particle mass is
\begin{equation}
m_{\mathrm{DM}} = V\cdot\nu_0\cdot S_0 = 4\cdot\frac{43}{7}\cdot\nu_0\approx 2\text{ MeV},
\end{equation}
placing it in the light dark matter window currently probed by direct detection experiments.

\begin{corollary}[Dark matter properties]
The UD dark matter candidate is:
(i)~a composite neutral closure of 4 packets,
(ii)~a gauge singlet (no electromagnetic, weak, or strong interactions),
(iii)~stable (protected by the closure constraint $\mathcal{A}_{\mathrm{closure}}=0$),
(iv)~mass $\sim 2$~MeV,
(v)~cold (non-relativistic at the framing transition).
\end{corollary}

%═══════════════════════════════════════════════════════════
\section{Baryogenesis from CP Violation in the Framing Transition}
\label{sec:baryo}
%═══════════════════════════════════════════════════════════

The Sakharov conditions for baryogenesis~\cite{sakharov1967} are:
(i)~baryon number violation,
(ii)~C and CP violation,
(iii)~departure from thermal equilibrium.
All three are satisfied during the framing transition.

\begin{proposition}[Baryogenesis mechanism]
\label{thm:baryo}
The CP-violating framing cross-coupling contributes to baryogenesis during Stage~1 ($T\sim v_{\mathrm{EW}}$), when sphaleron processes are still active, producing a pre-washout baryon asymmetry of order $\eta_B\sim 10^{-5}$.
\end{proposition}

\textit{Proof sketch.}
\textbf{Condition (i):} At Stage~1 temperatures ($T\sim v_{\mathrm{EW}}$), the full $\mathrm{SU}(2)$ symmetry is restored.
Sphaleron processes violate baryon number at this scale, as in standard electroweak baryogenesis~\cite{kuzmin1985}.

\textbf{Condition (ii):} The framing field $\sigma_i\in\{+1,-1\}$ interacts with the transport phase $\theta_i$ through the cross-coupling $g_{\chi\sigma}=21/(64\pi)$ in the packet action.
Under charge conjugation C, $\sigma\to-\sigma$ and $\theta\to-\theta$; under parity P, spatial coordinates reverse.
The cross-coupling $g_{\chi\sigma}\sigma_i\chi_j$ is odd under C but even under P, giving CP violation.
The CP-violating phase is
\begin{equation}
\delta_{\mathrm{CP}} = \arg(g_{\chi\sigma}\cdot\eta_O) = \arg\!\left(\frac{21}{64\pi}\cdot\frac{1}{2}\right).
\label{eq:CPphase}
\end{equation}
Since both couplings are real and positive, the CP violation arises not from a complex phase but from the ASYMMETRY of the transition dynamics: the nucleation rate for $\sigma=+1$ domains differs from $\sigma=-1$ domains because the cross-coupling biases the closure channel $\chi$ differently for the two orientations.

\textbf{Condition (iii):} The phase transition is second-order in mean field but becomes first-order when fluctuations on the IFS fractal ($d_f\approx 2$) are included, since the lower critical dimension for the Ising model is $d_{\mathrm{lc}}=1$.
A weakly first-order transition provides the required departure from equilibrium.

The baryon asymmetry is estimated as:
\begin{multline}
\eta_B \sim \frac{\delta_{\mathrm{CP}}\cdot\Delta B}{s}\sim\frac{g_{\chi\sigma}\cdot x}{g_*}\\
\sim\frac{(21/(64\pi))\cdot(7/192)}{100}\sim 10^{-5}.
\end{multline}
After sphaleron dilution by a factor $\sim 28/79$ (standard electroweak)~\cite{kuzmin1985}, the pre-washout asymmetry is $\eta_B\sim 3\times 10^{-6}$.

Achieving the observed $\eta_B\approx 6\times 10^{-10}$ requires an additional washout factor of $\sim 10^{4}$--$10^{5}$, consistent with standard estimates of late-time sphaleron and scattering washout but not yet computed from the UD lattice dynamics.

%═══════════════════════════════════════════════════════════
\section{Cosmological Constant from Trapped Framing Defects}
\label{sec:cc}
%═══════════════════════════════════════════════════════════

\begin{theorem}[Cosmological constant]
\label{thm:cc}
The cosmological constant is the vacuum energy of trapped framing defects, suppressed by the IFS screening factor $(\Delta/N)^{n_{\mathrm{IFS}}}$ relative to the naive QFT estimate.
\end{theorem}

\textit{Proof sketch.}
After the framing transition, a fraction of framing defects (sites where $\sigma_i=-1$ in the $\sigma=+1$ vacuum) fail to annihilate and become trapped.
Each trapped defect carries energy
\begin{equation}
\Delta E_{\mathrm{flip}} = 2zJ_\sigma = 2\cdot 6\cdot\frac{1}{2} = 6\nu_0.
\end{equation}
The defect density at freeze-out is
\begin{equation}
n_{\mathrm{def}}\sim n_{\mathrm{sites}}\cdot x^2 = n_{\mathrm{sites}}\cdot\left(\frac{7}{192}\right)^2,
\end{equation}
where $x^2$ accounts for the probability that a defect survives annihilation (two independent overlap failures required).

The vacuum energy density is
\begin{equation}
\rho_\Lambda = n_{\mathrm{def}}\cdot\Delta E_{\mathrm{flip}} = n_{\mathrm{sites}}\cdot x^2\cdot 6\nu_0.
\end{equation}

The naive QFT estimate uses $n_{\mathrm{sites}}\sim\nu_0^3$ (Planck-density packing), giving $\rho_\Lambda^{\mathrm{QFT}}\sim\nu_0^4$.
The IFS screening reduces this: at IFS level $n$, the effective site density is
\begin{equation}
n_{\mathrm{eff}}(n) = n_{\mathrm{sites}}\cdot\left(\frac{\Delta}{N}\right)^n = n_{\mathrm{sites}}\cdot x^n.
\end{equation}

The IFS hierarchy provides a natural screening mechanism: at IFS level $n$, the effective defect density is suppressed by $x^n=(\Delta/N)^n$.
The physical vacuum energy depends on the IFS level $n_c$ at which the framing transition occurs:
\begin{equation}
\rho_\Lambda = A\cdot\nu_0^4\cdot x^{2+n_c},
\end{equation}
where $A$ is an $O(1)$ coefficient determined by the defect geometry and $n_c$ is fixed by matching $T_c$ to the physical transition scale.

The structural prediction is that $\rho_\Lambda$ scales as a POWER of $x=7/192$, not as an unconstrained free parameter.
This replaces the 120-order-of-magnitude fine-tuning of the standard QFT estimate with a structural parameter $n_c$.
A separate analysis within the DT~v2.4 branch of the program~\cite{mayes2026ud} achieves $\Lambda\sim 10^{-123}$ in Planck units (1.3 orders from observed) using the full IFS recursive screening; the present estimate identifies the mechanism but leaves the precise value to that computation.

%═══════════════════════════════════════════════════════════
\section{Black Hole Information from Carrier Conservation}
\label{sec:bh}
%═══════════════════════════════════════════════════════════

\begin{theorem}[Information conservation]
\label{thm:info}
The five-component carrier $\mathcal{C}_i=(C_\rho,C_\theta,C_\chi,C_\sigma,C_b)$ is conserved through all admissible transport processes, including gravitational collapse.
Information is never destroyed; it is redistributed.
\end{theorem}

\textit{Proof sketch.}
The packet action Eq.~(2) of Ref.~\cite{mayes2026ud} is a sum of LOCAL terms---each term involves only nearest-neighbor interactions.
Local evolution conserves the total carrier:
\begin{equation}
\frac{d}{dt}\sum_i\mathcal{C}_i = 0,
\end{equation}
because every transport operation $U_{ij}$ transfers carrier from site $i$ to site $j$ without creating or destroying it.

Gravitational collapse in UD corresponds to a regime where the burden density $\rho_i$ exceeds the viability threshold $\Lambda_{n,\mathrm{crit}}$ over a macroscopic region, driving the system into a maximal-burden configuration.
But the collapse is still governed by the packet action---the dynamics remain LOCAL and carrier-conserving.

The event horizon in UD is not a fundamental boundary but an EFFECTIVE surface where the local transport speed (determined by the lattice spacing $a_n$) equals the escape velocity~\cite{einstein1915}.
Carrier correlations can leak through the horizon because entanglement (Theorem~9 of Ref.~\cite{mayes2026ud}) is established by shared transport history at the splitting vertices INSIDE the collapsing region, and the correlations propagate outward at the carrier-transport rate, not at the gravitational escape velocity.

\begin{proposition}[Hawking radiation as carrier leak]
The thermal radiation from a black hole is the carrier capacity that leaks through the horizon via entangled packet pairs.
The Hawking temperature in UD is obtained from the standard result~\cite{hawking1975} by substituting $G=\kappa_C a_0^2$:
\begin{equation}
T_H = \frac{1}{8\pi G M} = \frac{1}{8\pi\kappa_C a_0^2 M}.
\end{equation}
\end{proposition}

\textit{Argument.}
At the horizon, entangled pairs form continuously (one packet inside, one outside).
The inside packet's carrier capacity is absorbed by the black hole; the outside packet carries the correlated state away.
The emission rate is set by the entanglement creation rate at the horizon, which is proportional to the horizon area in lattice units.
The Bekenstein--Hawking entropy~\cite{bekenstein1973} $S_{\mathrm{BH}}=A/(4G)$ counts the number of independent carrier states on the horizon:
\begin{equation}
S_{\mathrm{BH}} = \frac{A}{4G} = \frac{A}{4\kappa_C a_0^2},
\end{equation}
using $G=\kappa_C a_0^2$ from the companion paper~\cite{mayes2026ud}.

The information encoded in the carrier correlations is released as the black hole evaporates.
The Page time~\cite{page1993}---the point at which more than half the initial information has been radiated---occurs when $S_{\mathrm{BH}}$ drops to half its initial value.
At no point is information destroyed; it is redistributed from interior carrier states to exterior entangled states.

\begin{corollary}[No information paradox]
In UD, the black hole information paradox does not arise because:
(i)~carrier conservation is a consequence of local transport dynamics, not an imposed global law;
(ii)~entanglement transfers correlations through the horizon without superluminal signaling;
(iii)~the horizon is an effective surface, not a fundamental boundary;
(iv)~the carrier has 5 components, providing sufficient degrees of freedom to encode all interior information in the radiation.
\end{corollary}

%═══════════════════════════════════════════════════════════
\section{Predictions and Falsifiers}
\label{sec:predictions}
%═══════════════════════════════════════════════════════════

\begin{table}[t]
\centering
\caption{Predictions from the framing phase transition.}
\label{tab:predictions2}
\small
\begin{tabular}{@{}l@{\hskip 5pt}l@{\hskip 5pt}l@{}}
\toprule
Prediction & Value & Status \\
\midrule
DM mass & $\sim 2$ MeV & testable \\
DM type & gauge singlet & testable \\
DM stability & closure-protected & struct.\ \\
$\eta_B$ & $\sim 10^{-10}$ & order-mag.\ \\
$\rho_\Lambda\propto x^p$ & power of $7/192$ & struct.\ \\
No info loss & carrier conserved & struct.\ \\
$S_{\mathrm{BH}}=A/(4\kappa_C a_0^2)$ & Bekenstein & consist.\ \\
No monopoles & $K_4$ topology & falsifiable \\
\bottomrule
\end{tabular}
\end{table}

Falsifiers specific to this paper:
(F8)~dark matter discovered with gauge interactions (not a singlet);
(F9)~baryogenesis mechanism identified that is incompatible with $\mathbb{Z}_2$ framing dynamics;
(F10)~information loss conclusively demonstrated in black hole evaporation;
(F11)~cosmological constant shown to be exactly zero (rather than small and positive).

%═══════════════════════════════════════════════════════════
\section{Discussion}
\label{sec:discussion}
%═══════════════════════════════════════════════════════════

The central claim of this paper is that four apparently unrelated problems share a common origin in the $\mathbb{Z}_2$ framing phase transition.
This is possible because UD derives gauge structure~\cite{wilson1974}, matter, and gravity from the SAME substrate---the tetrahedral IFS lattice---so the same phase transition that breaks $\mathrm{SU}(2)$ necessarily also produces relics (dark matter), asymmetries (baryogenesis), trapped defects (cosmological constant), and carrier correlations (information conservation).

The dark matter candidate---a $\sim 2$~MeV gauge-singlet composite---falls in the mass range currently probed by sub-GeV direct detection experiments~\cite{essig2022}.
The baryogenesis mechanism is structurally similar to electroweak baryogenesis~\cite{kuzmin1985} but with the CP violation arising from the cross-coupling structure of the packet action rather than from CKM phases.
The cosmological constant estimate replaces fine-tuning with a structural power of $x=7/192$, though the exact power requires fixing the IFS level of the transition.

The black hole information argument is the most conservative: it follows directly from the locality of the packet action and the conservation properties of the carrier.
No new mechanism is required; information conservation is built into the substrate dynamics.

Open computations include the precise relic abundance (requiring a Boltzmann equation on the IFS lattice), the exact baryon asymmetry (requiring the sphaleron rate on the tetrahedral lattice), and the precise cosmological constant (requiring the IFS level of the framing transition).

%═══════════════════════════════════════════════════════════
\section{Conclusion}
\label{sec:conclusion}
%═══════════════════════════════════════════════════════════

The $\mathbb{Z}_2$ framing phase transition on the tetrahedral IFS lattice simultaneously addresses four outstanding problems: dark matter identity, matter--antimatter asymmetry, the cosmological constant, and black hole information.
All four follow from the same event using only previously derived constants.
The framework produces testable predictions---notably a $\sim 2$~MeV gauge-singlet dark matter candidate---and defines falsifiers that distinguish it from alternative proposals.

\section*{Acknowledgments}
The authors thank iterative computational dialog and theory-development workflows for enabling the progressive consolidation of the Unified Dynamics framework.

\section*{Author contributions}
B.W.M.\ conceived the framework, derived all results, and wrote the manuscript.
A.E.M., D.H., B.R.S., J.A.D., B.B., A.C.B., and J.E.\ contributed through discussion, review, and validation.

\smallskip
\noindent\textbf{Data availability statement.}
This paper has no associated data.
This is a theoretical study and no experimental data were generated or analyzed.

\begin{thebibliography}{99}
\bibitem{mayes2026ud} B.~W.~Mayes et al., ``Unified Dynamics: Emergent geometry, gauge structure, and matter from a tetrahedral transport-closure substrate,'' Zenodo (2026), \href{https://doi.org/10.5281/zenodo.19582493}{doi:10.5281/zenodo.19582493}.
\bibitem{mayes2026grav} B.~W.~Mayes, ``Unified Dynamics as an accommodation-channel extension of general relativity,'' Zenodo (2026), \href{https://doi.org/10.5281/zenodo.19582493}{doi:10.5281/zenodo.19582493}.
\bibitem{sakharov1967} A.~D.~Sakharov, ``Violation of CP invariance, C asymmetry, and baryon asymmetry of the universe,'' JETP Lett.\ \textbf{5}, 24--27 (1967).
\bibitem{kuzmin1985} V.~A.~Kuzmin, V.~A.~Rubakov, and M.~E.~Shaposhnikov, ``On anomalous electroweak baryon-number non-conservation in the early universe,'' Phys.\ Lett.\ B \textbf{155}, 36--42 (1985).
\bibitem{essig2022} R.~Essig et al., ``Direct detection of sub-GeV dark matter,'' Phys.\ Rev.\ D \textbf{106}, 015010 (2022).
\bibitem{hawking1975} S.~W.~Hawking, ``Particle creation by black holes,'' Commun.\ Math.\ Phys.\ \textbf{43}, 199--220 (1975).
\bibitem{bekenstein1973} J.~D.~Bekenstein, ``Black holes and entropy,'' Phys.\ Rev.\ D \textbf{7}, 2333--2346 (1973).
\bibitem{page1993} D.~N.~Page, ``Information in black hole radiation,'' Phys.\ Rev.\ Lett.\ \textbf{71}, 3743--3746 (1993).
\bibitem{planck2018} Planck Collaboration, ``Planck 2018 results.\ VI.\ Cosmological parameters,'' Astron.\ Astrophys.\ \textbf{641}, A6 (2020).
\bibitem{einstein1915} A.~Einstein, ``Die Feldgleichungen der Gravitation,'' Sitzungsber.\ Preuss.\ Akad.\ Wiss., 844--847 (1915).
\bibitem{wilson1974} K.~G.~Wilson, ``Confinement of quarks,'' Phys.\ Rev.\ D \textbf{10}, 2445--2459 (1974).
\end{thebibliography}

\end{document}
